\documentclass[11pt]{article}
\usepackage{weekly}
\begin{document}
\weekheader{3}{13/06 \textendash{} 22/06/2026}{Agentic RAG, Tích hợp và Minh bạch hoá}{w3col}{Đã hoàn thành}

\section{Mục tiêu tuần}
Tuần thứ ba nâng hệ thống từ suy luận một lượt lên kiến trúc Agentic RAG: một tác tử (agent)
ReAct biết tự tra cứu corpus nội bộ, bổ sung tìm kiếm web khi cần và tổng hợp câu trả lời có
trích dẫn. Song song đó, nhóm tích hợp toàn bộ hệ thống lên lớp API và giao diện web, đồng
thời minh bạch hoá quá trình suy luận của agent cho người dùng cuối.

\section{Công việc đã thực hiện}

\paragraph{Giai đoạn 1 — Gom nguồn (Source Accumulation): }
Nhóm thiết kế tác tử ReAct vận hành theo vòng lặp Thought, Action, Observation. Agent được
trang bị hai công cụ: \texttt{docsearch} tra cứu corpus nội bộ (có tích hợp cổng lọc độ liên
quan như ở tuần 2) và \texttt{web\_search} tìm kiếm dự phòng trên Internet. Một policy ưu tiên
corpus được thiết lập qua system prompt: với mọi truy vấn, hành động đầu tiên luôn phải là
\texttt{docsearch}; agent chỉ được kích hoạt \texttt{web\_search} khi \texttt{docsearch} trả
về kết quả rỗng, hoặc khi corpus chỉ cung cấp thông tin dẫn chiếu chung chung mà truy vấn lại
cần thông tin thực địa cụ thể (địa chỉ, giờ làm việc, đường dây nóng theo địa bàn). Cơ chế này
kế thừa tư tưởng Corrective-RAG, được hiện thực bằng prompting kết hợp tín hiệu từ cổng lọc
thay vì fine-tune lại mô hình. Để tránh thiếu dữ kiện khi liệt kê hồ sơ, nhóm nâng giới hạn
ghép bằng chứng cho mỗi lượt từ khoảng 4000 lên 12000 token. Mọi đoạn corpus liên quan và mọi
kết quả web được đẩy vào một bể gom nguồn (Source Sink) dùng chung.

\paragraph{Công cụ tìm kiếm web (web\_search): }
Nhóm bọc Brave Search API thành một công cụ của agent. Công cụ lấy 5 kết quả đầu và tải nội
dung đầy đủ của 2 trang đầu tiên, bởi đoạn trích (snippet) thường thiếu các chi tiết thực địa
mà người dân cần như địa chỉ, giờ làm việc hay số điện thoại. Mỗi kết quả được bọc thành một
nguồn có thể trích dẫn, gán nhãn nguồn web (ký hiệu (web)) và gán điểm liên quan thấp một cách
có chủ ý do thông tin chưa qua thẩm định, để pha tổng hợp ưu tiên nguồn chính thống từ corpus.

\begin{figure}[H]
\centering
\resizebox{\textwidth}{!}{%
\begin{tikzpicture}
  \node[gbox] (q)       at (0,0)       {Truy vấn đầu vào\\(+ Contextualize)};
  \node[fbox] (thought) at (3.3,0)     {Thought\\(Suy luận)};
  \node[fbox] (action)  at (6.6,0)     {Action\\(Chọn công cụ)};
  \node[dbox] (doc)     at (9.9,1.35)  {docsearch\\(Corpus + Gating)};
  \node[abox] (web)     at (9.9,-1.35) {web\_search\\(Dự phòng)};
  \node[fbox] (obs)     at (13.2,0)    {Observation};
  \node[gbox] (sink)    at (16.5,0)    {Source Sink\\(Corpus + Web)};
  \node[fbox] (syn)     at (16.5,-3.7) {Giai đoạn 2: Tổng hợp\\có trích dẫn};
  \node[dbox] (ans)     at (12.3,-3.7) {Câu trả lời 【n】\\(Gắn nhãn nguồn Web)};
  \draw[farr] (q)--(thought); \draw[farr] (thought)--(action);
  \draw[farr] (action)--(doc); \draw[farr] (action)--(web);
  \draw[farr] (doc)--(obs); \draw[farr] (web)--(obs); \draw[farr] (obs)--(sink);
  \draw[farr] (obs.north) -- ++(0,1.85) -| (thought.north);
  \node[flab] at (4.9,2.5) {Loop $\leq N$};
  \draw[farr] (sink)--(syn); \draw[farr] (syn)--(ans);
  \begin{scope}[on background layer]
    \node[draw=black!35, dashed, rounded corners, fill=black!3,
          fit=(thought)(action)(doc)(web)(obs), inner sep=9pt] (p1) {};
  \end{scope}
  \node[font=\footnotesize] at (8.25,3.15) {GIAI ĐOẠN 1 — Source Accumulation (Vòng lặp ReAct)};
\end{tikzpicture}}
\caption{Tác tử ReAct hai giai đoạn: gom nguồn từ corpus (ưu tiên) và web (dự phòng), rồi
tổng hợp câu trả lời có trích dẫn.}
\label{fig:w3-react}
\end{figure}

\paragraph{Giai đoạn 2 — Tổng hợp (Synthesis): }
Tập nguồn đã gom được đưa qua đúng module tổng hợp có trích dẫn đã xây ở tuần 2, sinh câu trả
lời kèm các thẻ 【n】. Các trích dẫn trỏ về nguồn web được hệ thống tự động dán nhãn và đính
kèm hyperlink để phân biệt với nguồn chính thống, một cách xác định, không phụ thuộc việc LLM
có tuân thủ chỉ dẫn trong prompt hay không.

\paragraph{Hội thoại đa lượt, gợi ý và sơ đồ tư duy: }
Để hỗ trợ hội thoại đa lượt, nhóm bổ sung bước contextualize chạy trước Giai đoạn 1, giải các
tham chiếu chỉ định mơ hồ (thủ tục này, cơ quan ấy, hồ sơ đó) thành đối tượng cụ thể suy ra từ
lịch sử, tránh việc truy hồi bị lạc chủ đề. Hệ thống cũng sinh các câu hỏi gợi ý tiếp theo
(follow-up) dựa trên lịch sử hội thoại, và sinh sơ đồ tư duy (mind map) định dạng PlantUML cho
từng thủ tục, trong đó giấy tờ được gom thành nhóm bắt buộc chung và nhóm theo trường hợp.

\screenshotfig{img/week3-ui.png}{Toàn cảnh giao diện React: khung chat, sidebar hội thoại,
panel nguồn bên phải, ô gợi ý câu hỏi.}{Giao diện chính (React) của trợ lý hỏi đáp.}{fig:w3-ui}

\paragraph{Trực quan hoá suy luận (Reasoning Visualization): }
Trong suốt Giai đoạn 1, từng bước Thought, Action, Observation của agent được phát trực tiếp
lên giao diện theo thời gian thực thông qua Server-Sent Events. Mỗi bước hiển thị dưới dạng một
mục Bước N có thể bung hoặc gập, cho phép người dùng quan sát được lộ trình suy luận và tra cứu
của agent. Đây là yêu cầu thiết yếu về tính minh bạch (transparency) đối với các hệ thống AI
trong lĩnh vực hành chính công.

\screenshotfig{img/week3-reasoning.png}{Bảng các bước suy luận của agent hiển thị theo thời
gian thực: từng Bước N với Thought, Action, Observation (gọi docsearch rồi web\_search).}%
{Trực quan hoá quá trình suy luận của tác tử ReAct.}{fig:w3-reasoning}

\paragraph{Tích hợp Backend API và Frontend: }
Nhóm xây dựng một lớp backend FastAPI mỏng phơi ra các endpoint dưới dạng luồng SSE: gửi tin
nhắn, trả lời lại lượt cuối (regen), dừng luồng (stop), quản lý hội thoại, sinh gợi ý và đọc
ghi cấu hình. Mỗi luồng phát các loại sự kiện gồm câu trả lời tích luỹ, panel thông tin (các
bước suy luận và nguồn có tô sáng), danh sách nguồn có cấu trúc, biểu đồ trực quan trích dẫn và
sự kiện kết thúc kèm danh sách gợi ý. Giao diện được phát triển bằng Vite, React 19 và Tailwind,
gồm khung chat (render câu trả lời và liên kết trích dẫn), panel nguồn có tô sáng, bảng các
bước suy luận, ô gợi ý và các nút regen, stop.

\screenshotfig{img/week3-citation.png}{Câu trả lời có trích dẫn 【n】 và nhãn nguồn web; kèm
sơ đồ tư duy (mindmap) và các câu hỏi gợi ý tiếp theo.}{Trích dẫn nội dòng, nguồn web, mindmap
và gợi ý câu hỏi.}{fig:w3-citation}

\paragraph{Chiến lược thay thế ReWOO: }
Như một phương án cân bằng chi phí và độ trễ, nhóm tích hợp thêm kiến trúc ReWOO. Thay vì đan
xen suy luận và quan sát ở mỗi bước như ReAct, agent lập toàn bộ kế hoạch gọi công cụ từ
trước rồi mới thực thi và tổng hợp, giúp giảm số lượt gọi LLM. Cả ba chiến lược (Simple, ReAct,
ReWOO) dùng chung lõi truy hồi và lõi tổng hợp có trích dẫn, việc chuyển đổi chỉ là thiết lập
cấu hình lúc chạy.

\section{Thực nghiệm: đánh giá tầng truy hồi}
Truy hồi là nền tảng của tác tử, nên nhóm xây dựng một bộ dữ liệu đánh giá có nhãn và đo định
lượng tầng này.

\paragraph{Xây dựng bộ dữ liệu (Ground Truth): }
Bộ đánh giá gồm 210 câu: 180 câu sinh tự động bằng gpt-4o (phân tầng theo lĩnh vực) và 30 câu
viết tay gắn các thủ tục dân sinh có thật. 180 câu tự sinh chia đều thành sáu nhóm (mỗi nhóm
30 câu), mỗi nhóm kiểm một năng lực: khớp từ khoá (\texttt{factual}), diễn giải ngữ nghĩa
(\texttt{paraphrase}), câu tình huống gián tiếp (\texttt{scenario}), khía cạnh hồ sơ
(\texttt{aspect\_hoso}), khía cạnh phí/điều kiện (\texttt{aspect\_phi\_dk}) và truy vấn ngắn
(\texttt{keyword\_short}). Để chống thiên lệch BM25, ba nhóm \texttt{paraphrase}, \texttt{scenario},
\texttt{keyword\_short} chỉ cho mô hình thấy tên thủ tục (không cho text chunk), nên câu hỏi
không lặp nguyên văn từ khoá tài liệu. 30 câu viết tay gồm 20 câu in-scope và 10 câu ngoài
phạm vi để đo cổng lọc.

\paragraph{Độ đo và quy trình chấm: }
Khoá khớp gold là mã thủ tục. Với mỗi câu, hệ truy hồi trên toàn corpus, lấy top-20, khử trùng
theo thủ tục (giữ hạng tốt nhất), rồi tính Hit@$k$ (có ít nhất một gold trong top-$k$),
Recall@5, MRR (nghịch đảo hạng của gold đầu tiên) và nDCG@10. Với cỡ mẫu 180 câu, khoảng tin
cậy 95\% của tỉ lệ vào khoảng $\pm$4\%. Toàn bộ số liệu được sinh tự động bằng bộ script
\texttt{rag/eval/*}.

\paragraph{Kết quả: }
Bảng~\ref{tab:w3-retr} so sánh cấu hình lai (hybrid: vector + BM25, có rerank) với cấu hình chỉ
dùng vector. Hit@5 đạt 0,922 và Hit@10 đạt 0,972 --- tài liệu đúng gần như luôn nằm trong
top-10. Đáng chú ý, cấu hình chỉ-vector cho kết quả tương đương (Hit@5 0,922; MRR 0,781), cho
thấy embedding đa ngôn ngữ gánh gần hết còn BM25 đóng góp biên trên miền dữ liệu này. Điểm yếu
nằm ở hạng \#1 (Hit@1 0,661) với truy vấn ngắn (\texttt{keyword\_short} 0,533) và câu tình
huống (\texttt{scenario} 0,600), không phải ở recall. Về cổng lọc ``không tìm thấy'': tỉ lệ từ
chối đúng câu ngoài phạm vi đạt 100\% (10/10, không tài liệu nào lọt), nhưng recall in-scope
chỉ 0,65 --- cổng hơi gắt, đôi khi loại nhầm thủ tục đúng.

\begin{table}[H]
\centering\small
\renewcommand{\arraystretch}{1.2}
\begin{tabular}{@{}lrr@{}}
\toprule
Chỉ số & Hybrid & Chỉ-vector \\
\midrule
Hit@1 & 0,661 & 0,667 \\
Hit@3 & 0,861 & 0,889 \\
Hit@5 & 0,922 & 0,922 \\
Hit@10 & 0,972 & 0,956 \\
MRR & 0,775 & 0,781 \\
nDCG@10 & 0,823 & 0,824 \\
\bottomrule
\end{tabular}
\caption{Đánh giá tầng truy hồi (n=180): cấu hình lai so với chỉ-vector.}
\label{tab:w3-retr}
\end{table}

\section{Kết quả và sản phẩm bàn giao}
\begin{table}[H]
\centering\small
\begin{tabular}{@{}lp{8.7cm}@{}}
\toprule
Thành phần & Mô tả \\
\midrule
Reasoning ReAct & Tác tử hai giai đoạn, policy ưu tiên corpus, tổng hợp có trích dẫn. \\
\texttt{rag/agent\_tools.py} & Công cụ \texttt{web\_search} (Brave) làm dự phòng, có tải nội dung trang. \\
\texttt{app/api/} & Backend FastAPI phát SSE (chat, regen, stop, suggestions, settings). \\
\texttt{frontend/} & Giao diện React: chat, panel nguồn, bảng suy luận, gợi ý. \\
Tiện ích minh bạch & Trích dẫn 【n】 và tô sáng, nhãn nguồn web, mindmap, gợi ý follow-up. \\
\bottomrule
\end{tabular}
\end{table}

\section{Khó khăn và hướng xử lý}
\begin{itemize}
  \item Agent đôi khi tự phán đoán câu hỏi nằm ngoài phạm vi rồi bỏ qua \texttt{docsearch}:
        prompt ràng buộc luôn tra corpus trước, để dữ liệu quyết định thay vì agent đoán.
  \item Ngữ cảnh bị cắt làm thiếu thành phần hồ sơ: nhóm nâng giới hạn ghép bằng chứng từ
        khoảng 4000 lên 12000 token để giữ đủ các mục hồ sơ và lệ phí.
  \item Đoạn trích từ web thường thiếu chi tiết thực địa: công cụ tải nội dung trang đầu để
        pha tổng hợp có dữ liệu thật mà trích dẫn.
  \item Rò rỉ nguồn giữa các câu hỏi do công cụ dùng lại: nhóm cấp một bể gom nguồn mới và đặt
        lại cho mỗi yêu cầu.
  \item Lẫn cấp cơ quan (trung ương, tỉnh, xã) khi nêu nơi nộp hồ sơ: prompt buộc gán đúng địa
        chỉ theo cấp tương ứng và không suy diễn vượt nguồn.
\end{itemize}

\section{Checklist tuần 3}
\begin{itemize}[label={}]
  \cbDone Tác tử ReAct hai giai đoạn (ưu tiên corpus, dự phòng web, cổng lọc liên quan).
  \cbDone Công cụ \texttt{web\_search} (Brave) có tải nội dung trang.
  \cbDone Trích dẫn nội dòng và tô sáng, nhãn nguồn web, mindmap, gợi ý follow-up.
  \cbDone Contextualize hội thoại đa lượt; chiến lược thay thế ReWOO.
  \cbDone Tích hợp backend FastAPI (SSE) và frontend React.
  \cbDone Thực nghiệm đánh giá tầng truy hồi (bộ 180+30 câu; Hit@5 0,922, MRR 0,775).
\end{itemize}

\noindent Hướng tuần kế tiếp: bổ sung bộ nhớ hội thoại đa lượt (tóm tắt ngắn hạn và ký ức dài
hạn), tính năng định vị cơ quan xử lý thủ tục kèm bản đồ và chỉ đường, và đánh giá định lượng
toàn hệ thống (truy hồi, sinh câu trả lời và trích dẫn, bộ nhớ).

\end{document}
